 
\def\filename{leaflet-manual.tex}
\def\fileversion{v1.0d}   % change this when leaflet-manual changed, too.
\def\filedate{2012/06/04}
\def\docdate {2012/06/04} % change this when leaflet-manual changed, too.
\listfiles
\errorcontextlines=99
\documentclass[
notumble,
nofoldmark,
%%dvipdfm,
%%portrait,
%%titlepage,
%%nocombine,
%%a3paper,
%%debug,
%%nospecialtricks,
%%draft,
]{leaflet}


\renewcommand*\foldmarkrule{.3mm}
\renewcommand*\foldmarklength{5mm}

\usepackage[T1]{fontenc}
% \usepackage[T3]{inputenc}
\usepackage{textcomp}
\usepackage{mathptmx}
\usepackage{libertine}
\usepackage[scaled=0.9]{helvet}
\makeatletter
\def\ptmTeX{T\kern-.1667em\lower.5ex\hbox{E}\kern-.075emX\@}
\DeclareRobustCommand{\ptmLaTeX}{L\kern-.3em
        {\setbox0\hbox{T}%
         %\vb@xt@ % :-)
         \vbox to\ht0{\hbox{%
                            \csname S@\f@size\endcsname
                            \fontsize\sf@size\z@
                            \math@fontsfalse\selectfont
                            A}%
                      \vss}%
        }%
        \kern-.12em
        \ptmTeX}
\makeatother
\let\TeX=\ptmTeX
\let\LaTeX=\ptmLaTeX
\usepackage{shortvrb}
\MakeShortVerb{\|}
\usepackage{url}
\usepackage{graphicx}
\usepackage[dvipsnames,usenames]{color}
\definecolor{LIGHTGRAY}{gray}{.9}
\definecolor{lspblue1}{RGB}{208, 236, 246}
%%%%\renewcommand{\descfont}{\normalfont}
\newcommand\Lpack[1]{\textsf{#1}}
\newcommand\Lclass[1]{\textsf{#1}}
\newcommand\Lopt[1]{\texttt{#1}}
\newcommand\Lprog[1]{\textit{#1}}

\newcommand*\defaultmarker{\textsuperscript\textasteriskcentered}

\title{The document class \Lclass{leaflet}}
\author{%
  Rolf Niepraschk\\
  Walter Schmidt\\
  Hubert G\"a\ss lein}
\date{Last updated~\docdate\\printed \today}

\CutLine*{1}% Dotted line without scissors
\CutLine*{6}%   

% \AddToBackground{5}{%  Background of a small page
%   \put(0,0){\textcolor{Cerulean}{\rule{\paperwidth}{\paperheight}}}}
\AddToBackground*{1}{%  Background of a small page
  \put(0,0){\textcolor{lspblue1}{\rule{\paperwidth}{\paperheight}}}}
\AddToBackground*{2}{%  Background of a small page
  \put(0,0){\textcolor{lspblue1}{\rule{\paperwidth}{\paperheight}}}}

% \AddToBackground*{2}{% Background of a large page
%   \put(\LenToUnit{.5\paperwidth},\LenToUnit{.5\paperheight}){%
%     \makebox(0,0)[c]{%
%       \resizebox{.9\paperwidth}{!}{\rotatebox{35.26}{%
%         \textsf{\textbf{\textcolor{LIGHTGRAY}{BACKGROUND}}}}}}}}

\AddToBackground{2}{
    \put(0,580){
        \includegraphics[width=3\paperwidth]{lsp-rainbow-large.png}%
    } 
}
\AddToBackground{5}{
    \put(0,-100){
        \includegraphics[width=3\paperwidth]{lsp-rainbow-large.png}%
    } 
}

\AddToBackground{2}{
    \put(0,10){
        \includegraphics[width=.3\paperwidth]{lsp-rainbow-large.png}%
    } 
}

\AddToBackground{3}{
    \put(0,10){
        \includegraphics[width=.3\paperwidth]{lsp-rainbow-large.png}%
    } 
}

\AddToBackground{4}{
    \put(0,10){
        \includegraphics[width=.3\paperwidth]{lsp-rainbow-large.png}%
    } 
}


\begin{document}

% \maketitle
\thispagestyle{empty}

%%\LARGE

%%\tableofcontents
\noindent 
~
\parbox{.1\paperwidth}{
\vspace*{-.5cm}
\includegraphics[width=.3\paperwidth]{langsci_logo_nocolor.pdf} 
}
\parbox{.8\paperwidth}{
\raggedleft 
\vspace*{-2.5cm}{\includegraphics[width=.4\paperwidth]{lsp-rainbow-large.png}}
}

\vspace*{12cm}

\raggedleft{
\mbox{\Huge Language Science Press}
{\LARGE Open Access Monographs}}
\raggedright 

 \newpage

\section{ \sffamily   \Large About}
Vision [langfristig und pr�gnant. Nicht unbedingt direkt erreichbar]
Alle linguistischen Inhalte sind allen Menschen in einer nachhaltigen Struktur entgeltfrei zug�nglich. Die Zugeh�rigkeit zu einer Universit�t oder einem bestimmten Land darf f�r den Zugang keine Voraussetzung sein.
Weiterhin:

    Der Publikationsprozess entspricht den Anforderungen der Wissenschaftler

    Die Wahl des Publikationsortes wirkt sich nicht negativ auf die Karrierechancen aus

    Die  Verbreitung von Wissen erfolgt dezentral und ohne Beschr�nkungen

    Wissen wird nicht mehr monolithisch pr�sentiert, sondern ist granular, persistent und vernetzt verf�gbar 

    Nicht-textuelle Bestandteile (Video, quantitative Daten) sind selbstverst�ndliche Teile von Publikationen 

Mission [Wozu braucht man Language Science Press?]
Wir stellen eine  Plattform zur Verbreitung linguistischer Spitzenforschung zur Verf�gung, auf der Wissenschaftscommunities m�glichst autark ihre Forschungsergebnisse in Buchform als Open Access verbreiten k�nnen.
Ziele [konkret und messbar]

    Kundenzufriedenheit (Leser, Autoren, Herausgeber, Geldgeber)

    Schaffung einer starken, von Wissenschaftlern besessenen internationalen Marke, die f�r Qualit�t und Innovation steht

    Abdeckung aller Bereiche der Linguistik

    Stetiger Output

    Solide nachhaltige Struktur

 

\newpage 


\section{ \sffamily \Large Community}

\begin{itemize}


\item[$\rangle$]     Stefan M�ller, Prof. (Freie Universit�t Berlin)
\item[$\rangle$]     Matthias H�ning, Prof. Dr. (FU Berlin)
\item[$\rangle$]     Anatol Stefanowitsch, Prof. (Freie Universit�t Berlin)
\item[$\rangle$]     Alexander Koller, Prof. (Universit�t Potsdam)
\item[$\rangle$]     Hans C. Boas, Prof. (University of Texas at Austin)
\item[$\rangle$]     Frank Richter, PD Dr. (Eberhard-Karls-Universit�t T�bingen)
\item[$\rangle$]     Gert Webelhuth, Prof. (Goethe-Universit�t Frankfurt am Main)
\item[$\rangle$]     Gereon M�ller, Prof. (Universit�t Leipzig)
\item[$\rangle$]     Ivan Sag, Prof. (Stanford University)
\item[$\rangle$]     Tibor Kiss, Prof. (Ruhr-Universit�t Bochum)
\item[$\rangle$]     Martin Haspelmath, Prof. (MPI for Evolutionary Anthropology)
\item[$\rangle$]     Jonas Kuhn, Prof. (University Stuttgart)
\item[$\rangle$]     Manfred Sailer, Prof. (Goethe-Universit�t Frankfurt am Main)
\item[$\rangle$]     Emily M. Bender, Prof. (University of Washington)
\item[$\rangle$]     Friedemann Pulverm�ller, Prof. (Freie Universit�t Berlin)
\item[$\rangle$]     Frank Keller, Dr. (University of Edinburgh)
\item[$\rangle$]     Adele Goldberg, Prof. (Princeton University)
\item[$\rangle$]     Miriam Butt, Prof. (University of Konstanz)
\item[$\rangle$]     Laura Kallmeyer, Prof. (University of D�sseldorf)
\item[$\rangle$]     Shuly Wintner, Prof. (University of Haifa)
\item[$\rangle$]     Martine Grice, Prof. (University of Cologne)
\item[$\rangle$]     Uli Reich, Prof. (Freie Universit�t Berlin)
\item[$\rangle$]     Anke L�deling, Prof. (Humboldt-Universit�t zu Berlin)
\item[$\rangle$]     Artemis Alexiadou, Prof. (Universit�t Stuttgart)
\end{itemize}

 
\newpage 

\section{ \sffamily \Large Series}
 
\begin{itemize}
\item[$\rangle$] African Language Grammars and Dictionaries
\item[$\rangle$] Computational Models of Language Evolution
\item[$\rangle$] Conceptual Foundations of Language Science
\item[$\rangle$] Contemporary African Linguistics
\item[$\rangle$] Empirically Oriented Theoretical Morphology and Syntax
\item[$\rangle$] Implemented Grammars
\item[$\rangle$] Language Variation
\item[$\rangle$] Monographs on Comparative Niger-Congo
\item[$\rangle$] Studies in Caribbean Languages
\item[$\rangle$] Studies in Diversity Linguistics
\item[$\rangle$] Studies in Laboratory Phonology
\item[$\rangle$] Textbooks in Language Sciences
\item[$\rangle$] Topics at the Grammar-Discourse Interface
\item[$\rangle$] Translation and Multilingual Natural Language Processing
\end{itemize}

 

\newpage 

\section{ \sffamily \Large Books}
\begin{itemize}
 \item[$\rangle$]    The Alor-Pantar languages: History and typology	 
 \item[$\rangle$]    Prosodic detail in Neapolitan Italian	 
 \item[$\rangle$]    Natural causes of language: Frames, biases, and cultural transmission	 
 \item[$\rangle$]    A typology of marked-S languages	 
 \item[$\rangle$]    A grammar of Pite Saami	 
 \item[$\rangle$]     Adjective attribution	 
 \item[$\rangle$]     Syntax und Valenz  
 \item[$\rangle$]     Grammatical theory         
 \item[$\rangle$]     How mobile robots can self-organize a vocabulary	 
 \item[$\rangle$]     Language strategies for the domain of colour	 
 \item[$\rangle$]     Grammaire des constructions elliptiques  
 \item[$\rangle$]     The Talking Heads experiment: origins of words and meanings	  
 \item[$\rangle$]    The evolution of case grammar	 
 \item[$\rangle$]     The evolution of grounded spatial language	 
 \item[$\rangle$]     Einführung in die grammatische Beschreibung des Deutschen	 
 \item[$\rangle$]     A grammar of Yakkha	 
 \item[$\rangle$]     A grammar of Mauwake	  
\end{itemize}

   
    
   
   
   
   
\newpage 

\section{ \sffamily \Large Contact}

Press Contact\\
Mailing Address\\

Prof. Dr. Ste�fan M�l�ler\\
Language Science Press\\
Ha�bel�schwerd�ter Allee 45\\
14 195 Ber�lin\\
Germany\\

Principal Contact\\

Sebastian Nordhoff\\
Language Science Press\\
Ha�bel�schwerd�ter Allee 45\\
14 195 Ber�lin\\
Germany\\

Email: sebastian.nordhoff@langsci-press.org

http://www.langsci-press.org~

\vspace*{2cm}
% \framebox[3cm]{ \textbf{Mitmachen?}\\ Author\par Editor}

\begin{tabular}{lr}
\fbox{
  \parbox{3cm}{ 
    \textbf{Mitmachen?}\\ 
    \begin{itemize}
      \item[$\rangle$] author
      \item[$\rangle$] editor
      \item[$\rangle$] proofreader
      \item[$\rangle$] typesetter
      \item[$\rangle$] member
      \item[$\rangle$] supporter
      \item[$\rangle$] donate
    \end{itemize}
  }
}
&~
\raisebox{-2.2cm}{
\includegraphics[width=.4\textwidth]{qrcode.eps}
}
\end{tabular}


% \fbox{
% \textbf{Mitmachen?}
% \\
% Author
% \\
% Editor
% \\
% Proofreader
% \\
% Typesetter
% \\
% Membership
% \\
% Donate
% 
% }

% \begin{itemize}
%  \item[$\rangle$] autor
% \item[$\rangle$] hrsg
% \item[$\rangle$] proofreader
% \end{itemize}



\loggingall
\end{document}
\endinput
%%
%% End of file `leaflet-manual.tex'.
